% %%%%%%%%%%%%%%%%%%%%%%%%%%%%%%%%%%%%%%%%%%%%%%%%%%%%%%%%%%%%%%%%%%%%%
% %% SpDCache %%
% Generic Data-Cache
% %%%%%%%%%%%%%%%%%%%%%%%%%%%%%%%%%%%%%%%%%%%%%%%%%%%%%%%%%%%%%%%%%%%%%

\section{Generic Data-Caches}
\label{sec:spdc:gc}

\lettrine{M}{odern} processors and accelerators integrate data-caches to bridge the latency and bandwidth gap between computation units and main memory.
A cache is a small, fast memory that stores recently accessed data to anticipate future reuse.
Its efficiency relies on two main forms of locality:
\begin{itemize}
\item temporal locality, where recently accessed data is likely to be reused soon, and
\item spatial locality, where data stored near recently accessed addresses is likely to be used next.
\end{itemize}

\noindent
By exploiting these properties, caches reduce the average access latency and increase the overall bandwidth efficiency.

A cache is composed of several cache-lines, each grouping contiguous memory words.
The cache-line is the minimum granularity of data movement between memory levels.
When an access occurs, the cache controller searches for the requested address among the currently stored cache-lines.
If the address is found, a \gls{hit} occurs and the data is served directly to the core.
Otherwise, a \gls{miss} occurs, triggering a memory fetch from a lower cache level or from main memory.
If necessary, an existing cache-line is evicted to make room for the new data.
Each cache-line is associated with \glspl{tag} and validity bits that record its address and coherence state.

Several mapping policies exist.
In a direct-mapped cache, each memory block maps to a single cache-line, which is simple but prone to conflicts.
A $n$-\gls{way} \gls{set}-associative cache allows each block to be placed in one of $n$ possible cache-lines within a \gls{set}, improving \gls{hit} rate at moderate cost.
A \emph{fully associative} cache removes these constraints, as any block may occupy any cache-line, but requires more complex \gls{tag} searches.
When multiple candidates exist, a replacement policy such as \acrfull{lru} or \emph{Random} decides which entry is evicted.

Write policies also impact performance.
With a \emph{write-through} policy, every store operation updates both the cache and main memory, maintaining consistency but increasing traffic.
A \emph{write-back} policy delays the update until \gls{eviction}, reducing bandwidth usage at the cost of additional state bits (dirty flags) to maintain consistency.
For workloads dominated by \gls{seqAccesses}, these mechanisms are highly effective.
However, \glspl{irrAccess} tend to degrade cache efficiency by increasing \gls{miss} rates and provoking frequent \glspl{eviction}.

Modern processors implement hierarchical cache organizations to balance latency, capacity, and throughput.
A typical CPU integrates small L1 caches close to each core, larger L2 caches, and a high-capacity L3 cache shared at the chip level.
Accelerators and domain-specific architectures often employ local data storages or scratchpad memories serving the same purpose, providing low-latency access to frequently reused data (see Chapter~\ref{cha:soa}).

The generic cache assumes regular read and write operations with predictable locality.
Yet workloads operating on sparse or \gls{reduction}-intensive data violate this assumption.
In particular, ``random'' \glspl{reduction} are poorly supported, as conventional caches lack native mechanisms to merge updates efficiently or to retain relevant data in limited storage.

In the next sections, we introduce cache architectures supporting \gls{reduction} operations, specifically designed to handle \gls{irrData}.
These architectures generalize the concept of a data-cache to include \gls{reduction}-aware mechanisms, ensuring correctness while improving bandwidth efficiency.

In the following sections, we introduce cache architectures with native \gls{reduction} support, specifically designed to handle \gls{irrData}.
These architectures extend the classical cache model with \gls{reduction}-aware mechanisms, improving bandwidth efficiency under irregular access patterns.

\begin{highlight}{Conclusion}{Red}
    \textbf{Generic data-caches} efficiently exploit regular access patterns but \textbf{fail to address the irregularity} and accumulation requirements of sparse data computing.
    These limitations motivate the introduction of \textbf{\gls{reduction}-aware cache architectures}, capable of handling \gls{irrData}.
\end{highlight}